\chapter{Mod\`eles de diffusion et g\'en\'eration d'images}
\label{ch:diffusion}

\begin{quote}
\emph{<< Les mod\`eles de diffusion g\'en\`erent des images en apprenant \`a inverser un processus de bruitage. Partir d'un bruit pur, d\'ebruiter \'etape par \'etape, et une image coh\'erente \'emerge. >>}
\end{quote}

% --- hook ---
En 2020, Jonathan Ho et ses collègues de Berkeley redécouvrent une idée que la communauté avait largement abandonnée. Au lieu d'apprendre à générer directement une image à partir d'un bruit, on apprend à \emph{débruiter}. On part d'un bruit pur, et on enlève le bruit étape par étape, mille fois, jusqu'à ce qu'une image émerge. Stable Diffusion, DALL-E, Midjourney, Imagen --- toutes les images générées par IA que vous avez vues depuis 2022 reposent sur cette idée. Ce chapitre démonte la mécanique, code l'objectif d'entraînement DDPM, et fait tourner Stable Diffusion sur des prompts qui comptent.
\section{Le processus de diffusion}

\subsection{Processus avant (ajout de bruit)}

\'Etant donn\'ee une image $\mathbf{x}_0$, le processus avant ajoute du bruit gaussien sur $T$ pas de temps :
\[
q(\mathbf{x}_t | \mathbf{x}_{t-1}) = \mathcal{N}(\mathbf{x}_t; \sqrt{1-\beta_t}\,\mathbf{x}_{t-1},\; \beta_t \mathbf{I})
\]

o\`u $\beta_t$ est le sch\'ema de bruit. Apr\`es suffisamment d'\'etapes, $\mathbf{x}_T \approx \mathcal{N}(\mathbf{0}, \mathbf{I})$.

\subsection{Processus inverse (d\'ebruitage)}

Un r\'eseau de neurones apprend \`a inverser le bruit :
\[
p_\theta(\mathbf{x}_{t-1} | \mathbf{x}_t) = \mathcal{N}(\mathbf{x}_{t-1}; \boldsymbol{\mu}_\theta(\mathbf{x}_t, t),\; \sigma_t^2 \mathbf{I})
\]

Le mod\`ele pr\'edit le bruit $\boldsymbol{\epsilon}_\theta(\mathbf{x}_t, t)$ et l'utilise pour calculer $\boldsymbol{\mu}_\theta$.

\subsection{Objectif d'entra\^inement}

La perte DDPM simplifi\'ee est :
\[
L = \mathbb{E}_{t, \mathbf{x}_0, \boldsymbol{\epsilon}} \left[\| \boldsymbol{\epsilon} - \boldsymbol{\epsilon}_\theta(\mathbf{x}_t, t) \|^2 \right]
\]

\begin{minted}{python}
import torch
import torch.nn.functional as F

def diffusion_loss(model, x_0, noise_scheduler):
    """Simplified DDPM training step."""
    batch_size = x_0.shape[0]
    # \'Echantillonner des pas de temps al\'eatoires
    t = torch.randint(0, noise_scheduler.num_train_timesteps,
                      (batch_size,), device=x_0.device)
    # \'Echantillonner du bruit
    noise = torch.randn_like(x_0)
    # Ajouter du bruit aux images
    x_t = noise_scheduler.add_noise(x_0, noise, t)
    # Pr\'edire le bruit
    noise_pred = model(x_t, t).sample
    # Perte MSE
    loss = F.mse_loss(noise_pred, noise)
    return loss
\end{minted}

% ============================================================
\section{Sch\'emas de bruit}

Le sch\'ema de bruit $\{\beta_t\}_{t=1}^T$ contr\^ole \`a quelle vitesse le bruit est ajout\'e :

\begin{itemize}
  \item \textbf{Lin\'eaire :} $\beta_t$ cro\^it lin\'eairement de $\beta_1 = 10^{-4}$ \`a $\beta_T = 0.02$.
  \item \textbf{Cosinus :} ajout de bruit plus doux, meilleur pour les petites images.
  \item \textbf{Lin\'eaire mis \`a l'\'echelle :} utilis\'e par Stable Diffusion, r\'egl\'e pour l'espace latent.
\end{itemize}

\begin{minted}{python}
import numpy as np
import matplotlib.pyplot as plt

T = 1000
# Sch\'ema lin\'eaire
beta_linear = np.linspace(1e-4, 0.02, T)
alpha_linear = np.cumprod(1 - beta_linear)

# Sch\'ema cosinus
s = 0.008
steps = np.arange(T + 1) / T
f = np.cos((steps + s) / (1 + s) * np.pi / 2) ** 2
alpha_cosine = f[1:] / f[0]

plt.figure(figsize=(8, 4))
plt.plot(alpha_linear, label="Linear")
plt.plot(alpha_cosine, label="Cosine")
plt.xlabel("Timestep")
plt.ylabel("Cumulative alpha (signal remaining)")
plt.title("Noise schedules")
plt.legend()
plt.tight_layout()
plt.savefig("noise_schedules.png", dpi=150)
plt.show()
\end{minted}

% ============================================================
\section{L'architecture U-Net}

Les mod\`eles de diffusion utilisent typiquement un \textbf{U-Net} : un encodeur-d\'ecodeur avec connexions r\'esiduelles. Le U-Net prend une image bruit\'ee et le pas de temps, et pr\'edit le bruit.

Composants cl\'es :
\begin{itemize}
  \item \textbf{Blocs de downsampling} : convolutions qui r\'eduisent la r\'esolution spatiale.
  \item \textbf{Blocs d'upsampling} : convolutions transpos\'ees qui augmentent la r\'esolution.
  \item \textbf{Skip connections} : concat\`enent les caract\'eristiques de l'encodeur aux caract\'eristiques du d\'ecodeur.
  \item \textbf{Plongement du pas de temps} : plongement sinuso\"idal de $t$, inject\'e dans chaque bloc.
  \item \textbf{Cross-attention} : pour la g\'en\'eration conditionn\'ee par texte, le U-Net attribue son attention aux plongements de texte.
\end{itemize}

% ============================================================
\section{Stable Diffusion}

Stable Diffusion op\`ere dans l'\emph{espace latent} : les images sont compress\'ees par un encodeur VAE, la diffusion se passe dans l'espace latent, et le d\'ecodeur VAE reconstruit l'image.

\begin{minted}{python}
from diffusers import StableDiffusionPipeline
import torch

pipe = StableDiffusionPipeline.from_pretrained(
    "stabilityai/stable-diffusion-2-1-base",
    torch_dtype=torch.float16,
)
pipe = pipe.to("cuda")

prompt = "A serene African village at sunset, digital art, highly detailed"
image = pipe(
    prompt,
    num_inference_steps=30,
    guidance_scale=7.5,
).images[0]

image.save("african_village.png")
image
\end{minted}

\begin{aitip}
\texttt{guidance\_scale} contr\^ole \`a quel point l'image suit le prompt. Des valeurs de 7--9 marchent bien pour la plupart des prompts. Des valeurs plus hautes produisent des interpr\'etations plus litt\'erales mais peuvent r\'eduire la qualit\'e d'image.
\end{aitip}

% ============================================================
\section{Prompts n\'egatifs et param\`etres}

\begin{minted}{python}
image = pipe(
    prompt="Portrait of a scientist in a laboratory, photorealistic",
    negative_prompt="blurry, low quality, deformed, cartoon",
    num_inference_steps=50,
    guidance_scale=8.0,
    height=512,
    width=512,
).images[0]
image.save("scientist.png")
\end{minted}

% ============================================================
\section{G\'en\'eration image-vers-image}

Partir d'une image existante et la modifier :

\begin{minted}{python}
from diffusers import StableDiffusionImg2ImgPipeline
from PIL import Image

pipe_img2img = StableDiffusionImg2ImgPipeline.from_pretrained(
    "stabilityai/stable-diffusion-2-1-base",
    torch_dtype=torch.float16,
).to("cuda")

init_image = Image.open("sketch.png").resize((512, 512))

result = pipe_img2img(
    prompt="A beautiful watercolor painting of a landscape",
    image=init_image,
    strength=0.75,  # 0=pas de changement, 1=r\'eg\'en\'eration compl\`ete
    guidance_scale=7.5,
).images[0]
result.save("watercolor.png")
\end{minted}

% ============================================================
\section{ControlNet : g\'en\'eration guid\'ee}

ControlNet ajoute un conditionnement spatial (contours, profondeur, pose) \`a Stable Diffusion :

\begin{minted}{python}
from diffusers import StableDiffusionControlNetPipeline, ControlNetModel
from controlnet_aux import CannyDetector
from PIL import Image
import torch

# Charger ControlNet pour les contours Canny
controlnet = ControlNetModel.from_pretrained(
    "lllyasviel/sd-controlnet-canny", torch_dtype=torch.float16
)
pipe_cn = StableDiffusionControlNetPipeline.from_pretrained(
    "stabilityai/stable-diffusion-2-1-base",
    controlnet=controlnet,
    torch_dtype=torch.float16,
).to("cuda")

# Extraire les contours d'une image de r\'ef\'erence
canny = CannyDetector()
reference = Image.open("building_photo.png")
edges = canny(reference, low_threshold=100, high_threshold=200)

# G\'en\'erer avec guidage par contours
result = pipe_cn(
    prompt="A futuristic building, cyberpunk style, neon lights",
    image=edges,
    num_inference_steps=30,
).images[0]
result.save("cyberpunk_building.png")
\end{minted}

% ============================================================
\section{G\'en\'eration par lots et contr\^ole de seed}

\begin{minted}{python}
import torch

# G\'en\'eration reproductible avec seeds
generator = torch.Generator("cuda").manual_seed(42)

images = pipe(
    prompt=["A cat astronaut", "A dog scientist", "A bird musician"],
    num_inference_steps=30,
    guidance_scale=7.5,
    generator=generator,
).images

for i, img in enumerate(images):
    img.save(f"generated_{i}.png")
\end{minted}

\begin{exercise}
\begin{enumerate}
  \item G\'en\'erez 5 images avec le m\^eme prompt mais des seeds diff\'erentes. Quelle variation observez-vous ?
  \item Exp\'erimentez avec les valeurs de \texttt{guidance\_scale} : 1, 5, 7.5, 12, 20. Documentez l'effet sur la qualit\'e d'image et le respect du prompt.
  \item Utilisez image-vers-image pour transformer un croquis simple en illustration d\'etaill\'ee. Essayez diff\'erentes valeurs de \texttt{strength}.
  \item Appliquez ControlNet avec d\'etection de contours Canny pour g\'en\'erer des variantes architecturales d'une photo de b\^atiment.
  \item G\'en\'erez une image, puis utilisez img2img pour cr\'eer 3 variations avec \texttt{strength} croissant (0.3, 0.5, 0.8).
\end{enumerate}
\end{exercise}

\begin{ethicsbox}
La g\'en\'eration d'images soul\`eve des questions \'ethiques s\'erieuses :
\begin{itemize}
  \item \textbf{Deepfakes :} g\'en\'erer des images r\'ealistes de personnes r\'eelles sans consentement.
  \item \textbf{Droits d'auteur :} mod\`eles entra\^in\'es sur de l'art prot\'eg\'e sans autorisation des artistes.
  \item \textbf{Biais :} les mod\`eles peuvent perp\'etuer des st\'er\'eotypes (par ex., g\'en\'erer surtout des visages \`a peau claire pour << professionnel >>).
  \item \textbf{Contenus nocifs :} les mod\`eles peuvent g\'en\'erer des images violentes ou explicites.
\end{itemize}
Utilisez toujours les v\'erificateurs de s\^uret\'e (activ\'es par d\'efaut dans diffusers), filigranez les images g\'en\'er\'ees par IA, et ne g\'en\'erez jamais d'images non consenties de personnes r\'eelles.
\end{ethicsbox}

% ============================================================
\section{R\'esum\'e du chapitre}

\begin{itemize}
  \item Les mod\`eles de diffusion apprennent \`a inverser un processus de bruitage : partir du bruit, d\'ebruiter \'etape par \'etape.
  \item L'objectif d'entra\^inement est simple : pr\'edire le bruit ajout\'e.
  \item Les sch\'emas de bruit (lin\'eaire, cosinus) contr\^olent la vitesse de diffusion.
  \item U-Net avec cross-attention est l'architecture standard pour la g\'en\'eration conditionn\'ee par texte.
  \item Stable Diffusion travaille dans l'espace latent pour l'efficacit\'e.
  \item ControlNet ajoute un guidage spatial (contours, profondeur, pose) au processus de g\'en\'eration.
  \item Guidance scale, prompts n\'egatifs et nombre d'\'etapes sont les param\`etres cl\'es de g\'en\'eration.
  \item Un d\'eploiement \'ethique demande filtres de s\^uret\'e, filigranage, et consentement.
\end{itemize}
